% Turkish (Türkçe) -- what the reading editions' preamble leaves to the language.
% Input by lib/frank-preamble.tex as \input{\FrankLib/lang/\BookLang.tex};
% docs/languages.md section 6 lists the hooks every file here must define and
% section 12 the ones it may, and the manual's "Adding a language" chapter
% says in one line each what happens when one is missing.
%
% A Latin-script file like it.tex: no font of its own, no marks to strip, two
% passes.  What is Turkish about it is the vocabulary line -- an agglutinative
% word is a clause, so the gloss is a segmentation -- and the three verb forms
% that segmentation needs, which is why the second label below is not the
% default and the first, pres., does not mean the Persian present stem.

% ---- the babel language and the switches ----------------------------------
% \FrankLangName is what \foreignlanguage and \begin{otherlanguage} are given
% everywhere in the core preamble: babel must know this name, or every passage
% comes out in the roman with a "Unknown language" error.
\newcommand{\FrankLangName}{turkish}
\FrankRTLfalse           % chunk left, gloss right; \pw needs no bidi isolate
\FrankHasBarefalse       % Turkish spelling leaves nothing out, so pass 3 would repeat pass 1
\FrankHasAltfalse        % one hand only: no nastaliq, no gothic
\FrankHasReadingfalse    % the spelling is the reading
\FrankHasVertfalse       % horizontal, like every Latin script

% ---- fonts ------------------------------------------------------------------
% No face is NAMED here: \FrankPickFont walks the registry's tex.main and then
% tex.main_fallbacks and takes the first one installed, so the .tex and the
% .json cannot drift apart (docs/languages.md section 12).  Turkish asks for
% none of that -- tex.main is null, so the text is set in the roman the core
% preamble already loaded (TeX Gyre Pagella, which has ı İ ğ ş and the dotted
% capital), the same face as the glosses.  What import=tr is here for is the
% Turkish locale: the case pairs, so that an \MakeUppercase anywhere in the
% chrome gives i -> İ and I -> ı and not the English pair, and the hyphenation
% patterns of hyph-tr, which install.sh --pdf installs (hyphen-turkish) and
% checks.  A machine that skipped it makes babel say "Hyphen rules for
% 'turkish' set to \l@english" as an Info and then break Turkish at English
% points for the rest of the run -- biliy-ordu, gülümsedi whole and unbroken,
% where the patterns give bi-li-yordu and gü-lüm-sedi.  Agglutination makes
% Turkish the language this costs most: its words are long enough that nearly
% every one of them reaches a line end sooner or later.
\babelprovide[import=tr]{turkish}

% ---- the vocabulary line -----------------------------------------------------
% \vb{infinitive}{}{present in -iyor}{}{aorist}{}{meaning} -- the -mek/-mak
% citation form, then the two forms a Turkish gloss cannot do without, both
% in the bare third person singular.  The aorist is the one form the
% infinitive does not predict -- gelir but bakar, alır but olur -- so it is
% the Turkish counterpart of Persian's present stem.  The -iyor present is
% the form a reader meets most, and the one where a verb is hardest to
% recognise: demek gives diyor, yemek yiyor, aramak arıyor, oynamak oynuyor,
% gitmek gidiyor.
%
% THE SECOND SLOT USED TO BE THE STEM, written as the endings "really attach"
% to it -- gid-, ed-, tad- -- and that was wrong for half the paradigm: before
% a consonant the stem is git- (gitti, gitmiş, the imperative git!), so a
% reader told "stem gid-" built *gidti.  And the stem is the infinitive less
% -mek/-mak in every other verb, so the slot taught nothing that the page did
% not already print.  The -iyor present carries what it tried to carry (the
% d of gidiyor) and the narrowing of de-, ye- and every a/e stem besides; the
% t/d of gider was already in the aorist.
%
% All three sound slots stay empty: the spelling already says how each form
% is said.  The one extra, in ONE parenthesis after the meaning and only
% where the gloss-language verb would not lead the reader to it, is the case
% the verb's complement takes: -e, -i, -den, -de, ile --
%   \vb{bakmak}{}{bakıyor}{}{bakar}{}{to look at (-e)}
% and the form the text actually has, segmented in \textit, still comes last
% after a semicolon (docs/lang/tr.md).  That file names these same three in
% this same order, and must be changed with these two lines.
% The reader takes the same pair from the registry (lib/languages.json,
% "vb_labels": ["pres.", "aor."], the three forms in words as "vb_forms"); the
% two must agree, or the PDF and the reader of one book label the same form
% differently.  lib/newlang.py --check compares them.
\newcommand{\FrankVbPres}{pres.}
\newcommand{\FrankVbPast}{aor.}

% ---- the two Lua filters ---------------------------------------------------
% frank_strip(s): the bare form of a chunk -- what pass 3 shows.  Written from
% the registry's `strip`; the identity when there is nothing to take off.
% frank_latin(s): the label's digits as ASCII, because a PDF destination name
% cannot carry the language's own figures.  Written from `digits`.
% Both are the identity for Turkish: the alphabet writes every vowel it has,
% so nothing is added for pass 1 to carry and dropped for pass 3, and the
% figures are already 0-9.  They still have to exist -- the core calls them
% for every language.  No #, % or literal backslash in a \directlua body: TeX
% reads it first (NOTES section 9 of the Persian book).
\directlua{
  function frank_strip(s) return s end
  function frank_latin(s) return s end
}

% ---- the "How to read this" page -----------------------------------------------
% Printed by lib/frank-frontmatter.tex, before the first chapter: the only
% place the reader is told what the passes are for, and the only place the
% hyphens of the vocabulary line are explained -- without that paragraph
% ev-ler-imiz-den is a word with hyphens in it and nothing more.
% \pw sets a word as Turkish inside the English and \textit sets a sound or a
% segmentation, which is the division every gloss line makes (\dw is \pw then
% \textit): the reader meets the two typographies here in the roles they keep
% for the rest of the book.  \pw is defined by the core preamble, after this
% file is read, and is only expanded when the page is set.
\newcommand{\FrankHowTo}{%
Every passage comes twice.

\smallskip
\textbf{First} the whole sentence, and nothing else. Try to read it. Get as
far as you can and do not worry about what you cannot get — the point is
the attempt, made before any help arrives.

\smallskip
\textbf{Second} the same sentence in small chunks, Turkish on the left, and
on the right three lines: how it is pronounced, then how each word is built,
then what it means. Now you find out how much you had right. Turkish is
written as it is said, one letter to one sound, so there is no third pass:
nothing can be taken away.

\smallskip
The middle line takes the words apart, because a Turkish word is often a
whole clause: \pw{evlerimizden} is \textit{ev-ler-imiz-den}, house + plural +
our + from. The hyphens are seams: put the pieces back and you have the word
as it is spelled; the names run in that order. A stem that changed
shape is followed by the form you would look up: \textit{kitab-ı} is
\pw{kitap}.

\smallskip
Verbs are given as the infinitive, the present in \textit{-iyor} and the
aorist, both as \textit{he} or \textit{she} would say them: \pw{gelmek},
\pw{geliyor}, \pw{gelir}. The present is the form you will meet most and
the one that hides the verb best — \pw{demek} is \pw{diyor}, \pw{aramak}
\pw{arıyor}. The aorist vowel is the one thing a Turkish verb will not tell
you — \pw{gelir} but \pw{bakar}, \pw{alır} but \pw{olur}. A case ending in
brackets after the meaning is the one the verb wants on what it governs:
\pw{bakmak} (\textit{-e}), to look \emph{at}, takes \pw{pencereye}, at the
window. In a compound with \pw{etmek} or \pw{olmak}, only the noun's meaning
is given.

\smallskip
The pronunciation line is given only where the spelling misleads, which is
seldom. A macron is a long vowel — an Arabic or Persian loan printed short
(\pw{kâr} \textit{kār} profit, \pw{kar} snow), or the vowel \pw{ğ} has
melted into (\pw{dağ} \textit{dā}); between front vowels \pw{ğ} is a
\textit{y} instead (\pw{değil} \textit{deyil}). \textit{ä} is the open
\textit{e} of \pw{gel} and \pw{ben}. \textit{ˈ} goes before the stressed
syllable, and only where the stress is not final, where Turkish otherwise
puts it: \textit{ˈyalnız}, and always before the negative \pw{-ma},
\pw{-me} — \textit{ˈgelme} is ``do not come'', \textit{gelˈme}
``coming''.
}
